\documentclass{article}

\usepackage{tikz} % before corl_2026, so its colours are defined with xcolor
\usetikzlibrary{arrows.meta,positioning,fit,shapes.geometric,calc,backgrounds}
\usepackage[preprint,nonatbib]{corl_2026}
% as in ../main.tex: references in citation order, the 30 evaluations' primary sources first,
% so reference number = evaluation id
\usepackage[square,numbers]{natbib}
\bibliographystyle{unsrtnat}
\usepackage{graphicx}
\usepackage{booktabs}
\usepackage{amsmath,amssymb}
\usepackage{longtable}
\usepackage{listings}
\usepackage{placeins}

% The extended report: everything in the paper (the shared sections under ../sections/),
% every alternative figure, the ablations, the elicitation details and the final ensemble.
% Numbers come from ../generated/macros.tex (and ../generated/noctx/macros.tex when the
% no-context ablation has a run), written by scripts/regen.sh.
\graphicspath{{../generated/}}
\input{../generated/macros.tex}
\newcommand{\reportdir}{../}
% Notation and macros shared by the paper (main.tex) and the extended report (extended/main.tex).
\newcommand{\EVSI}{\mathrm{EVSI}}
\newcommand{\EVPI}{\mathrm{EVPI}}
\newcommand{\pist}{\pi^\star}
\newcommand{\EVSIst}{\EVSI^{*}}
\newcommand{\etast}{\eta^{*}}
\newcommand{\PS}{P_{S}}
\newcommand{\Cb}{C_\mathrm{b}}
\newcommand{\Cr}{C_\mathrm{r}}
\newcommand{\anonlink}{\url{https://anonymous.4open.science/r/eval-voi-E8B2/}}
% \reportdir: path from the compiling document to this directory. \ifextended: true in the
% extended report only (it sets \extendedtrue); the paper never mentions that report.
\providecommand{\reportdir}{}
\newif\ifextended

% as in ../main.tex: drop the template's first-page enlargement, which lands on page 2
\makeatletter
\renewcommand{\@notice}{\@float{noticebox}[b]\footnotesize\@noticestring\end@float}
\makeatother
\extendedtrue
\newif\ifnoctx
\IfFileExists{../generated/noctx/macros.tex}{\input{../generated/noctx/macros.tex}\noctxtrue}{}

\lstset{basicstyle=\ttfamily\scriptsize, breaklines=true, columns=fullflexible, keepspaces=true,
  frame=single, framerule=0.3pt, xleftmargin=2pt, xrightmargin=2pt}

% "left out of the paper because ...": one line under each figure that is not in the paper
\newcommand{\why}[1]{\FloatBarrier\par\smallskip\noindent\emph{Not in the paper because} #1\par}

\title{Do physical-AI safety evaluations pay for themselves? Extended report}

\author{
  Anonymous Author(s)\\
  Affiliation\\
  \texttt{email}
}

\begin{document}
\voiNociteEvaluations % the evaluations are references 1 to \voiNScenarios, in id order
\maketitle

% Shared by main.tex and extended/main.tex. Lines marked % CHECK-DATA depend on the data.
\begin{abstract}
A safety evaluation is a measurement, worth money when its result can change a release decision and so avoid harm.
Physical-AI safety evaluations look expensive next to a question set for a large language model (LLM).
We price each evaluation by the expected value of sample information (EVSI) for the release decision it informs, per dollar of build and run cost.
The parameters of this decision model are elicited as probability distributions from an ensemble of \voiNMembers{} LLMs for \voiNPhysical{} physical-AI and \voiNLLM{} LLM safety evaluations, and Monte Carlo sampling carries their uncertainty.
Among the evaluations tested, a randomly chosen physical-AI evaluation gives an undecided developer more value per dollar than a randomly chosen LLM evaluation only about \voiEtaStarPSQMed{} of the time, though the best physical-AI evaluation places in the top half of all evaluations on \voiEtaStarPBestPhysicalTopHalf{} of draws. % CHECK-DATA
The groups differ in the stakes of the decision, about \voiStakesRatio{} times higher for an LLM release, and the gap would close as the stakes of physical-AI releases grow with every robot deployed. % CHECK-DATA
In both groups, the evaluations whose result can change the decision repay their build within a few runs. % CHECK-DATA
\end{abstract}


\keywords{Safety evaluation, value of information, physical AI}

\noindent This report carries the paper's text unchanged (Sections~\ref{sec:intro} to~\ref{sec:conclusion}), then every figure and analysis that did not fit in four pages, the ablations, the elicitation details, the evaluations and their sources, and the ensemble.

% Shared by main.tex and extended/main.tex. Lines marked % CHECK-DATA depend on the data.
\begin{figure}[t]
  \centering
  \includegraphics[width=\linewidth]{fig_headline.pdf}
  \caption{Value of one run against the cost $C$ to build and run it, every parameter at its median; labels are evaluation ids (reference numbers). (a) $\EVSIst$, the value to an undecided decision maker, all \voiNScenarios{} evaluations. (b) EVSI at the elicited prior belief that the hazard is present: \voiNZeroPhysical/\voiNPhysical{} physical-AI and \voiNZeroLLM/\voiNLLM{} LLM evaluations have zero value of information, the prior being so decisive that no result would change the decision. Colour: risk domain; triangles: physical AI. Dashed lines: constant value per dollar; up and left is better. % CHECK-DATA
  }
  \label{fig:headline}
\end{figure}


% Shared by main.tex and extended/main.tex. Lines marked % CHECK-DATA depend on the data.
\section{Introduction}
\label{sec:intro}

Developers of large language models (LLMs) run safety evaluations before a release and report them in system cards~\citep{shevlane2023extreme,anthropic2026rsp,openai2025preparedness}; a virology question set~\citep{gotting2025vct} runs in minutes through a programming interface. Physical-AI safety evaluations test a robot, a vehicle or a model that controls one, from question answering about physical danger~\citep{jindal2025danger} to world-model rollouts~\citep{geminirobotics2025veo}, jailbreaks of real robots~\citep{robey2024robopair} and vehicle track protocols~\citep{iihs2024paeb}. Many need hardware, scenes and staff, so they look too expensive to be worth running.

Cost alone cannot say whether an evaluation is worth running, because an evaluation is a measurement, worth what its result does. A safety evaluation is worth the harm its result avoids by improving the release decision, so a cheap test whose result cannot move the decision is worth nothing and an expensive test that would stop a harmful release is worth much. Decision analysis quantifies this as the expected value of sample information (EVSI)~\citep{raiffa1961applied,howard1966information}, used in health economics to choose which studies to fund~\citep{claxton1999irrelevance}. AI safety evaluations deserve the same scrutiny: their budget is finite and should go to the evaluations that give the most safety per dollar.

We estimate the value of information of \voiNPhysical{} physical-AI safety evaluations and compare it with that of \voiNLLM{} LLM safety evaluations from system cards, each per dollar of build and run cost. The parameters of the decision model are elicited as probability distributions from an ensemble of \voiNMembers{} LLMs, and Monte Carlo draws carry their uncertainty into rankings and group comparisons. We report rankings and comparisons rather than absolute values, because elicited parameters are imperfect: a bias the elicitors share across all evaluations moves both groups alike.

% Shared by main.tex and extended/main.tex. Lines marked % CHECK-DATA depend on the data.
\section{Method}
\label{sec:method}

\begin{figure}[t]
  \centering
  \input{\reportdir pipeline.tex}
  \caption{Method. The candidate evaluations (left) each get a decision model (centre): the hidden state $\theta$ with prior $p$ drives the result $x$ through sensitivity $s$ and specificity $t$; the developer sees $x$ and chooses the action $a$; the utility $U$ depends on $\theta$ and $a$ through $B$ and $K$. Two independent prompts elicit the decision parameters (green) and the instrument parameters (purple). Monte Carlo propagates the elicited uncertainty to the value per dollar (right).}
  \label{fig:method}
\end{figure}

\emph{Decision model.} We use a deliberately simple model that gives order-of-magnitude values and rankings, not point predictions (Fig.~\ref{fig:method}). One developer faces one release decision: respond (delay and mitigate) or deploy as planned. A hidden binary state $\theta$ is 1 when the hazardous property is present, so that responding is correct; its prior is $p=P(\theta{=}1)$. One run of the evaluation returns a flag $x\in\{0,1\}$ with sensitivity $s=P(x{=}1\mid\theta{=}1)$ and specificity $t=P(x{=}0\mid\theta{=}0)$. Responding gains $B$ if $\theta=1$ and loses $K$ if $\theta=0$. We value $B$ and $K$ for society, harm avoided and benefits forgone, and assume the developer acts on that value, as published safety frameworks state; the developer's own payoff is an ablation. With stakes $\Lambda=B+K$ and threshold $\pist=K/\Lambda$, acting on a belief $\pi$ is worth $V(\pi)=\Lambda\,[\pi-\pist]^+$ more than deploying, where $[z]^+=\max(z,0)$. A flag arrives with probability $P_1=ps+(1-p)(1-t)$ and moves the belief to $\pi_1=ps/P_1$; a pass moves it to $\pi_0=p(1-s)/(1-P_1)$. The value of one run and its efficiency are
\begin{equation}
\EVSI = P_1\,V(\pi_1)+(1-P_1)\,V(\pi_0)-V(p), \qquad \eta=\EVSI/C, \qquad C=C_\mathrm{build}+C_\mathrm{run},
\label{eq:evsi}
\end{equation}
where $C_\mathrm{build}$ is the cost to build the evaluation from scratch and $C_\mathrm{run}$ the cost of one run against one system. EVSI is zero exactly when no result moves the belief across the threshold: $\pi_0\ge\pist$, so the developer responds regardless, or $\pi_1\le\pist$, so it deploys regardless~\citep{pauker1980threshold}. This happens when the developer is already nearly certain ($p$ far from $\pist$ relative to the evaluation's power), and always when the evaluation is uninformative ($s+t=1$). At fixed stakes, EVSI is largest over thresholds when the developer is exactly undecided ($\pist=p$), where it equals the indifference value $\EVSI^\circ=\Lambda\,p(1-p)\,|s+t-1|$: stakes times the variance of the state times Youden's index $J=s+t-1$~\citep{youden1950index}. A built evaluation informs $n$ decisions over its life (model releases, product versions). Building pays after $n^*=C_\mathrm{build}/(\EVSI-C_\mathrm{run})$ decisions, and never if one run costs more than it is worth.

\emph{Evaluations.} The LLM safety evaluations are those named in LLM system cards. We rank them by how often system cards name them, whether the instrument is public, how checkable the score is and whether one run gives one result, and include those above a fixed score (\inextended{Section~\ref{sec:catalogue}}). They span cyber, biological, loss-of-control, manipulation and societal harms. The physical-AI safety evaluations are published evaluations of embodied systems that score the rate of a physically harmful event. Each has a public score, a cost or validity fact and one build-and-run instrument. They span fidelity levels from text question answering (level 0) to track testing (level 8). Each evaluation gets one release decision written from its sources, a definition of $\theta$, and two paragraphs of sourced facts: one about the decision, one about the evaluation. We invent no number in the model; every number comes from the elicitation.

\emph{Elicitation.} An ensemble of language models stands in for the expert panel that a study of this scale cannot afford. Language models have been used to elicit priors~\citep{capstick2024autoelicit,selby2025hadenough}, following the practice of expert elicitation~\citep{ohagan2006uncertain}, and a diverse LLM crowd forecasts about as well as a human crowd~\citep{schoenegger2024wisdom}. Each evaluation gets two independent prompts. The decision prompt asks for $p$, $B$ and $K$ and never describes the evaluation. The instrument prompt asks for $s$, $t$, $C_\mathrm{build}$, $C_\mathrm{run}$ and $n$ and never states the stakes. For each parameter the member writes its reasoning, then a value it would be surprised to see the truth fall below (5th percentile) or above (95th), then its median. The members are \voiNMembers{} Claude models (\voiMemberList), each answering each prompt twice. % CHECK-DATA: members and repeats of the analysed run
Answers are validated (strict JSON, ordered percentiles, $s>1-t$) and retried once. Of \voiAttempts{} attempts, \voiValidShare{} were valid; \voiInvalidRefusal{} were refusals, all on one biosecurity prompt, and \voiInvalidJson{} were unparseable. % CHECK-DATA The language models are the measurement instrument of this study; this is our main use of generative AI.

\emph{Monte Carlo.} Each elicitation's three percentiles are fitted by least squares with a Beta distribution ($p$, $s$, $t$) or a lognormal ($B$, $K$, $C_\mathrm{build}$, $C_\mathrm{run}$, $n$). The pooled belief per parameter is the equal-weight mixture of the fitted distributions, the linear opinion pool~\citep{stone1961opinionpool,clemen1999combining}. We draw \voiDraws{} parameter vectors per evaluation, with parameters independent, and rank all evaluations on each draw. The central estimate evaluates Eq.~\eqref{eq:evsi} at the median over elicitations of each elicited median.

\emph{Comparing the groups.} The headline metric is $A=P(\eta_\mathrm{phys}>\eta_\mathrm{LLM})$, the probability that a randomly chosen physical-AI evaluation has a higher efficiency than a randomly chosen LLM evaluation, with ties counted as one half. It equals the area under the receiver operating characteristic curve and the Mann-Whitney $U$ statistic divided by the number of pairs~\citep{mcgraw1992cles,vargha2000stochastic}. The percentile curve generalises it to the $q$-th percentile LLM evaluation. We compute both at the central estimate and on every draw. An exact Mann-Whitney test on the central estimates, with p-value $p_\mathrm{MW}$, is the frequentist companion.

% Shared by main.tex and extended/main.tex. Lines marked % CHECK-DATA depend on the data.
\section{Results}
\label{sec:results}

\emph{Placement.} At the central estimate a random physical-AI evaluation is worth more per dollar than a random LLM evaluation with probability $A=\voiEtaACentral$, and the exact Mann-Whitney test separates the groups ($p_\mathrm{MW}=\voiEtaMWp$). % CHECK-DATA
On the indifference value, which does not depend on the threshold, every LLM evaluation ranks above every physical-AI one ($A=\voiEtaIndACentral$, $p_\mathrm{MW}=\voiEtaIndMWp$). Run-only efficiency $\EVSI/C_\mathrm{run}$, which ignores the build cost, gives $A=\voiEtaRunACentral$. % CHECK-DATA
Over Monte Carlo draws $A$ sits near one half (median \voiEtaAQMed, 90\% interval \voiEtaAQLo{} to \voiEtaAQHi). On most draws many evaluations of both groups are worth exactly zero, and ties count one half, so the draw-level $A$ mostly measures how often neither result can change the decision; the central estimate and the indifference value carry the comparison. % CHECK-DATA
The five highest efficiencies belong to LLM evaluations, led by \voiTopShort{} ($\eta=\voiTopEta$). The best physical-AI evaluations are the robot-command refusal suite ($\eta=\voiEtaScLLMRobotCommandRefusalSuites$) and ASIMOV agentic orchestration ($\eta=\voiEtaScASIMOVAgenticSafetyOrchestration$); the pedestrian braking track protocol reaches $\eta=\voiEtaScPedestrianAEBTrackProtocols$. % CHECK-DATA

\emph{Why: the stakes, not the cost.} The groups look alike on what an evaluation controls. The median sensitivity is \voiMedSPhysical{} for physical-AI and \voiMedSLLM{} for LLM evaluations, the median specificity \voiMedTPhysical{} and \voiMedTLLM, the median build cost \voiMedCBuildPhysical{} and \voiMedCBuildLLM, the median run cost \voiMedCRunPhysical{} and \voiMedCRunLLM, and the median reuse count \voiMedNPhysical{} and \voiMedNLLM. They differ in the decision. The median stakes $B+K$ are \voiMedStakesLLM{} for an LLM release and \voiMedStakesPhysical{} for a physical-AI one, and the median stakes per evaluation dollar are \voiMedStakesPerCostLLM{} against \voiMedStakesPerCostPhysical, a factor of \voiStakesPerCostRatio. The median indifference efficiency differs by a similar factor (\voiMedEtaIndLLM{} against \voiMedEtaIndPhysical). % CHECK-DATA
Physical-AI priors are also higher (median $p$ \voiMedPPhysical{} against \voiMedPLLM), so more of them are cases where the developer responds regardless. The result can change the decision for \voiChangesShareLLM{} of the LLM and \voiChangesSharePhysical{} of the physical-AI evaluations; \voiNRespondsRegardlessPhysical{} of the \voiNPhysical{} physical-AI evaluations cannot, because the developer responds regardless. There, the useful step is to mitigate first and evaluate the mitigated system. % CHECK-DATA

\emph{Scale.} EVSI and the indifference value grow in proportion to the stakes at a fixed threshold $K/(B+K)$. Physical-AI evaluations would draw level with LLM ones on the indifference value at about \voiEtaIndRatio{} times today's stakes per decision. % CHECK-DATA

\emph{Break-even reuse.} Where one run is worth more than it costs to run, building pays fast: the median break-even reuse count at the central estimate is \voiNstarMedianPhysical{} decisions for physical-AI and \voiNstarMedianLLM{} for LLM evaluations, below one, against an elicited reuse count near \voiMedNPhysical. Across all evaluations of a group, the median probability that the elicited reuse reaches the break-even count is \voiPPaysMedianPhysical{} (physical AI) and \voiPPaysMedianLLM{} (LLM). % CHECK-DATA

\emph{Robustness and fidelity.} Averaged over evaluations, the uncertainty in $\eta$ correlates most with the uncertainty in $K$ (mean absolute Spearman \voiRhoK), and least with the cost parameters (\voiRhoCBuild{} for $C_\mathrm{build}$). The costs, the parameters the intuition is about, barely move the value. Across the \voiNLevel{} physical-AI evaluations, the Spearman correlation of fidelity level with the elicited Youden's index is $\voiLevelRhoJ$ and with cost $\voiLevelRhoC$: in this set, higher fidelity buys neither a more discriminating test nor a costlier one. % CHECK-DATA

% Shared by main.tex and extended/main.tex. Lines marked % CHECK-DATA depend on the data.
\begin{figure}[t]
  \centering
  \includegraphics[width=0.96\linewidth]{fig_rows.pdf}
  \caption{Per evaluation over the draws, rows sorted by the median rank, coloured by risk domain; physical-AI rows outlined and bold. Left: rank by $\etast$ (1 = best): median (tick), interquartile range (thick) and 90\% interval (thin). Right: break-even run count $n^*=\Cb/(\EVSI-\Cr)$ on the draws where one run is worth more than its run cost (violin, log scale; margin: that share of draws). % CHECK-DATA
  }
  \label{fig:rows}
\end{figure}

% Shared by main.tex and extended/main.tex. Lines marked % CHECK-DATA depend on the data.
\section{Limitations and conclusion}
\label{sec:limitations}

Every parameter is an LLM's belief. The \voiNMembers{} members come from \voiNMembers{} developers, none of them the developer of the fact-extracting model, and all read the same sourced facts; still, the parameters are not validated against human experts, LLM intervals are often too narrow~\citep{hobor2026bayesian}, and no evaluation publishes a sensitivity and specificity against real harm. The extracted facts are a second weak point: an agent chose which sentences of the sources to show.\ifextended{} \ifdefined\voinoctxEtaStarPSQMed With no facts shown at all, $\PS$ by $\etast$ is \voinoctxEtaStarPSQMed{} (against \voiEtaStarPSQMed{} with facts) and the stakes ratio \voinoctxStakesRatio{} (against \voiStakesRatio) (Section~\ref{sec:noctx}).\fi\fi{} % CHECK-DATA
Each decision is written from the evaluation's own sources, so the comparison is between the decisions as much as between the evaluations; pricing both groups against one shared decision would separate the two. Published physical-AI safety evaluations are few and uneven, so the result describes today's evaluations, not physical AI as such.

\label{sec:conclusion}
Under our model, physical-AI safety evaluations give an undecided developer less value per dollar than LLM safety evaluations because far less is at stake per release. In both groups, the evaluations that can change a decision repay their build quickly. Pricing evaluations by their value of information gives a budget rule that takes measured sensitivities, specificities and costs as they become available. Code and data: \anonlink. % CHECK-DATA


%===============================================================================
\FloatBarrier
\section{Further figures}
\label{sec:alternatives}

Every figure below is written by the same analysis run as the paper's figures. Evaluation ids are listed in Table~\ref{tab:scenarios}.

\begin{figure}[!htbp]
  \centering
  \includegraphics[width=\linewidth]{fig_rank_star.pdf}
  \caption{Rank by $\etast$ per evaluation (1 = best) over the Monte Carlo draws, the left panel of Fig.~\ref{fig:rows} as intervals: median (tick), interquartile range (thick bar) and 90\% interval (thin bar), rows sorted by the median rank, coloured by risk domain. Median interquartile range: \voiEtaStarRankIQRMedian{} places (physical AI \voiEtaStarRankIQRMedianPhysical, LLM \voiEtaStarRankIQRMedianLLM).}
  \label{fig:rankstar}
\end{figure}

\begin{figure}[!htbp]
  \centering
  \includegraphics[width=\linewidth]{fig_rank_eta.pdf}
  \caption{The same by $\eta$, where the elicited prior counts. On a draw where an evaluation's prior is already decisive it is worth zero and shares the average rank of the tied zeros, so a rank moves between its positive-value place and the shared place of the zeros. Median interquartile range: \voiEtaRankIQRMedian{} places (physical AI \voiEtaRankIQRMedianPhysical, LLM \voiEtaRankIQRMedianLLM).}
  \label{fig:ranketa}
\end{figure}
\why{the paper ranks by $\etast$, which has no ties; this figure shows what the elicited prior does to the ranking.}

\begin{figure}[!htbp]
  \centering
  \includegraphics[width=\linewidth]{fig_breakeven.pdf}
  \caption{Break-even run count $n^*=\Cb/(\EVSI-\Cr)$ per evaluation on the draws where one run is worth more than its run cost (violin, log scale; tick: median), rows sorted by the median $n^*$; right margin: the share of draws on which $n^*$ is finite. An $n^*$ below 1 means the first run already repays the build.}
  \label{fig:breakeven}
\end{figure}
\why{the paper shows it beside the ranks (Fig.~\ref{fig:rows}, right); here the rows are sorted by $n^*$ itself.}

\begin{figure}[!htbp]
  \centering
  \includegraphics[width=\linewidth]{fig_best_physical_rank.pdf}
  \caption{On each draw, the rank of the best-ranked physical-AI evaluation among all \voiNScenarios{} evaluations; the curve is the share of draws on which that rank is at most $N$, by $\etast$ (solid) and by $\eta$ (dotted). By $\etast$: top 1 \voiEtaStarPBestPhysicalTopOne, top 3 \voiEtaStarPBestPhysicalTopThree, top 5 \voiEtaStarPBestPhysicalTopFive, top 10 \voiEtaStarPBestPhysicalTopTen; with every parameter at its median the best physical-AI evaluation (\voiEtaStarBestPhysicalShort) ranks \voiEtaStarBestPhysicalRankCentral.}
  \label{fig:bestphys}
\end{figure}
\why{the paper quotes its top-5 and top-10 values; the curve itself is a detail.}

\begin{figure}[!htbp]
  \centering
  \includegraphics[width=0.62\linewidth]{fig_curve.pdf}
  \caption{The probability that a randomly chosen physical-AI evaluation has a higher $\eta$ than the $q$-th percentile LLM evaluation, against $q$: median and 90\% band over draws, and with every parameter at its median (dashed). At $q=50$: \voiCurveCentralFifty{} at the median parameters, \voiCurveMedFifty{} over draws.}
  \label{fig:curve}
\end{figure}
\why{it generalises $\PS$ to any percentile, but the zeros under $\eta$ make the band hard to read; the paper reports $\PS$ by $\etast$ instead.}

\begin{figure}[!htbp]
  \centering
  \includegraphics[width=\linewidth]{fig_percentile_violins.pdf}
  \caption{Per physical-AI evaluation, the distribution over draws of its percentile among the LLM evaluations by $\eta$ (0: below every LLM evaluation, 100: above every one). Tick: median over draws; diamond: every parameter at its median. The percentile can only take one more distinct value than there are LLM evaluations, plus the ties, so the waves in the violins are that discreteness smoothed by the kernel, not structure in the elicitations.}
  \label{fig:violins}
\end{figure}
\why{it answers the per-evaluation question ("which LLM evaluation is this one as good as?") but the discrete percentiles make it look like more than it is.}

\begin{figure}[!htbp]
  \centering
  \includegraphics[width=\linewidth]{fig_roc.pdf}
  \caption{Receiver operating characteristic (ROC) curve of ``physical AI'' given the score, for $\eta$ (left) and $\etast$ (right). A threshold sweeps down the score; each point gives the share of physical-AI evaluations (y) and of LLM evaluations (x) above it. Black: every parameter at its median; orange: median and 5--95\% band over draws, pointwise at a fixed grid of x values. Evaluations tied across the groups enter at one threshold, so the curve runs diagonally there; its area is $\PS$, ties counted as one half (\voiEtaPS{} and \voiEtaStarPS{} at the median parameters). Below the diagonal, physical-AI evaluations score lower.}
  \label{fig:roc}
\end{figure}
\why{it shows where along the score the groups separate, which the paper has no room for.}

\begin{figure}[!htbp]
  \centering
  \includegraphics[width=\linewidth]{fig_params.pdf}
  \caption{The elicited inputs. One panel per parameter; one point per valid answer (its median), coloured by member; black tick: the median across members, the value the single estimate uses. USD for $B$, $K$, $\Cb$ and $\Cr$.}
  \label{fig:params}
\end{figure}
\why{it makes every input transparent, which answers a reviewer who asks what the members said and where they disagree.}

\begin{figure}[!htbp]
  \centering
  \includegraphics[width=\linewidth]{fig_sensitivity.pdf}
  \caption{Spearman correlation of each parameter's Monte Carlo draws with $\etast$, per evaluation: a global sensitivity measure in the manner of partial rank correlation~\citep{marino2008sensitivity}. Large magnitudes mark the inputs whose uncertainty drives that evaluation's value per dollar. Mean absolute value over evaluations: $p$ \voiRhoIndP, $s$ \voiRhoIndS, $t$ \voiRhoIndT, $B$ \voiRhoIndB, $K$ \voiRhoIndK, $\Cb$ \voiRhoIndCBuild, $\Cr$ \voiRhoIndCRun. The same against $\eta$: $p$ \voiRhoP, $s$ \voiRhoS, $t$ \voiRhoT, $B$ \voiRhoB, $K$ \voiRhoK, $\Cb$ \voiRhoCBuild, $\Cr$ \voiRhoCRun{} (figure fig\_sensitivity\_eta in the repository).}
  \label{fig:sensitivity}
\end{figure}
\why{the paper gives the averages in one sentence.}

\begin{figure}[!htbp]
  \centering
  \includegraphics[width=\linewidth]{fig_level.pdf}
  \caption{Physical-AI evaluations by fidelity level (\voiLevelMin: text-only planning, \voiLevelMax: track testing), every parameter at its median. Top: $\eta$ (left) and $\etast$ (right); open markers on the bottom row are evaluations worth zero because the prior is already decisive. Spearman correlation with level: \voiLevelRhoEta{} ($p=\voiLevelRhoEtap$) for $\eta$ and \voiLevelRhoEtaStar{} ($p=\voiLevelRhoEtaStarp$) for $\etast$, over \voiNLevel{} evaluations. Bottom, secondary: Youden's index $J=s+t-1$ (\voiLevelRhoJ, $p=\voiLevelRhoJp$) and cost $C$ (\voiLevelRhoC, $p=\voiLevelRhoCp$).}
  \label{fig:level}
\end{figure}
\why{the paper gives the correlation in one sentence; a robotics reader may want the points.}

\begin{figure}[!htbp]
  \centering
  \includegraphics[width=\linewidth]{fig_members.pdf}
  \caption{Cross-member agreement. Per parameter, each member's median per evaluation (x) against the median of the other members (y). Points on the diagonal agree. Mean pairwise Spearman correlation between members: $p$ \voiMemberRhoP, $s$ \voiMemberRhoS, $t$ \voiMemberRhoT, $B$ \voiMemberRhoB, $K$ \voiMemberRhoK, $\Cb$ \voiMemberRhoCBuild, $\Cr$ \voiMemberRhoCRun.}
  \label{fig:members}
\end{figure}
\why{it is evidence about the inputs, not a result.}

\begin{table}[!htbp]
  \centering
  \caption{Per-evaluation numbers with every parameter at its median, ordered by $\etast$. Rank: by $\etast$, with its 90\% interval over draws. $P(\mathrm{rank}\le k)$: share of draws on which the evaluation ranks in the top $k$ by $\etast$. $P(\EVSI>0)$: share of draws on which the result can change the decision. $n^*$: break-even run count.}
  \label{tab:headline}
  \resizebox{\linewidth}{!}{\input{../generated/tab_headline.tex}}
\end{table}

%===============================================================================
\FloatBarrier
\section{Ablations}
\label{sec:ablations}

\emph{No curated facts.}
\label{sec:noctx}
The main run's prompts show each member two paragraphs of sourced facts. The no-context ablation (protocol \texttt{final\_noctx}) renders the same prompts to the same members with no facts and no facts heading: the decision prompt names only the agent, the decision and $\theta$; the instrument prompt only the evaluation's title, agent, decision, $\theta$ and instrument.
\ifnoctx
Without facts, a random physical-AI evaluation beats a random LLM evaluation by $\etast$ with probability \voinoctxEtaStarPSQMed{} over the draws (90\% interval \voinoctxEtaStarPSQLo{} to \voinoctxEtaStarPSQHi; exact Mann-Whitney $p=\voinoctxEtaStarMWp$), against \voiEtaStarPSQMed{} with facts; by $\eta$, \voinoctxEtaPSQMed{} against \voiEtaPSQMed. The median stakes ratio is \voinoctxStakesRatio{} against \voiStakesRatio; the share of evaluations whose prior is already decisive is \voinoctxZeroSharePhysical{} (physical AI) and \voinoctxZeroShareLLM{} (LLM) against \voiZeroSharePhysical{} and \voiZeroShareLLM; the best physical-AI evaluation ranks \voinoctxEtaStarBestPhysicalRankCentral{} against \voiEtaStarBestPhysicalRankCentral. Of \voinoctxAttempts{} answers, \voinoctxValidShare{} were valid. % CHECK-DATA
\begin{figure}[!htbp]
  \centering
  \includegraphics[width=\linewidth]{noctx/fig_headline.pdf}
  \caption{Figure~\ref{fig:headline} without the curated facts.}
  \label{fig:noctx}
\end{figure}
\else
This ablation has no stored run yet; this paragraph fills in when it has.
\fi

\emph{Run-only value per dollar.} A developer who adopts a published evaluation pays only the run. Its value per dollar is $\eta_\mathrm{run}=\EVSI/\Cr$. Under it, $\PS=\voiEtaRunPSQMed$ over the draws (\voiEtaRunPSQLo{} to \voiEtaRunPSQHi; exact Mann-Whitney $p=\voiEtaRunMWp$ at the median parameters), against \voiEtaPSQMed{} under $\eta$. % CHECK-DATA

\emph{$\eta$ against $\etast$.} Under $\eta$ the elicited prior counts and \voiNZeroPhysical{} physical-AI and \voiNZeroLLM{} LLM evaluations are worth zero at the median parameters (ids \voiZeroIdsPhysical; \voiZeroIdsLLM); among the evaluations whose result can change the decision, $\PS=\voiEtaPSChanging$ over \voiEtaPSChangingN{} pairs. Under $\etast$ no evaluation is worth zero and $\PS=\voiEtaStarPS$ at the median parameters. % CHECK-DATA

%===============================================================================
\FloatBarrier
\section{Elicitation details}
\label{sec:elicitation}

\emph{Prompts.} Each evaluation has two prompts, rendered from the two templates below. The decision prompt fills in the agent, the decision, the definition of $\theta$ and the decision facts; it never describes or names the evaluation (it has no title field), so the stakes and the prior cannot lean on the evaluation's framing. The instrument prompt fills in the evaluation's title, agent, decision, definition of $\theta$, the instrument (one sentence naming the evaluation and what it is run against) and the instrument facts; it never states $p$, $B$ or $K$. The two are separate calls because in the pilots a single prompt that described the evaluation shifted the elicited prior $p$ towards the evaluation's own framing. The facts are sentences taken from the cited sources, each traceable to a cached copy of its source text.

\emph{Where the facts come from.} The decision and instrument facts were written and checked by language-model agents, not by hand (research/scenarios\_draft/PROVENANCE.md in the repository, with every prompt verbatim). The model was Claude Fable 5.1, run as Claude Code subagents that could fetch web pages and read the repository, before any elicitation existed. For each evaluation, one agent fetched the primary sources, cached their full text and wrote a cited decision paragraph and a cited instrument paragraph of 120 to 250 words each. A separate agent checked every sentence against the cached text and the separation rule (no instrument facts in the decision paragraph, no stakes in the instrument paragraph), and a third fixed what it found. A script then removes the citation markers and writes the scenarios the prompts render. This guarantees that every sentence an elicitor reads traces to a cached source and was checked by a second agent. It does not guarantee a neutral selection: an agent chose which facts to include within the word limit, which is why the ensemble has no Anthropic model (Section~\ref{sec:final}) and why the no-context ablation exists (Section~\ref{sec:noctx}).

\emph{Order and format.} Per parameter, the member writes a reasoning paragraph of three to six sentences (the reasoning is a summary the member commits to, not its full chain of thought, which the members produce internally at the same reasoning effort), then the 5th and 95th percentiles as surprise questions (``a value you would be surprised to see the true value fall below or above''), then the median. Extremes first is standard practice against anchoring on the best guess~\citep{ohagan2006uncertain}. Numbers are plain decimals; the parser also reads scientific notation.

\emph{Validation.} An answer is valid when it is strict JSON with exactly the parameters of its prompt, each with $p_5<p_{50}<p_{95}$; probabilities lie in $(0,1)$; USD amounts are positive; the median $s$ exceeds one minus the median $t$ (the evaluation is informative); and the prior's median lies in $[0.001,0.999]$. A failed attempt is retried once. Failures are classed as unparseable JSON, wrong schema, a violated constraint, a failed fit, a refusal (no JSON and the text declines, or a content filter), a truncated answer, or a transport error.

\emph{Fitting and combining.} Each answer's three percentiles are fitted by least squares on the quantiles: a Beta distribution for $p$, $s$ and $t$, a lognormal for $B$, $K$, $\Cb$ and $\Cr$. Two parameters fitted to three quantiles leave a residual, stored as a fit-quality value. The combined belief per parameter is the equal-weight mixture of the fitted distributions over every valid answer: on each draw one member is chosen at random and its distribution sampled (the linear opinion pool~\citep{stone1961opinionpool,clemen1999combining}). A single distribution fitted to all members' quantiles would be forced to one mode and one width; the mixture keeps disagreement visible. The seven parameters are drawn independently, and each member's component independently per evaluation, so a member's systematic view is not carried through a draw; this narrows the intervals in a direction we cannot sign.

\emph{Members and health.} Table~\ref{tab:health} gives attempts, the valid share, the failures by class and the cost per member. Of \voiAttempts{} answers, \voiValid{} (\voiValidShare) were valid, \voiInvalidRefusal{} were refused, \voiInvalidJson{} were unparseable and \voiInvalidConstraint{} violated a constraint; the run cost \voiUSD. The refusals were Gemini's content filter on the decision prompts of two CBRN evaluations (VCT, ABC-Bench) and the unparseable answer GLM's; each of these slots was filled on the retry, so every slot has a valid answer. % CHECK-DATA

\begin{table}[!htbp]
  \centering
  \caption{Elicitation health per member: attempts, share valid after the one retry, failures by class (unparseable JSON, wrong schema, violated constraint, failed fit, refusal, truncated answer, and the transport errors cli, http, api, provider and other), cost in USD.}
  \label{tab:health}
  \resizebox{\linewidth}{!}{\input{../generated/tab_health.tex}}
\end{table}

\emph{Decision template.}
\lstinputlisting{../../templates/decision.md}

\emph{Instrument template.}
\lstinputlisting{../../templates/instrument.md}

%===============================================================================
\FloatBarrier
\section{The evaluations}
\label{sec:catalogue}

Table~\ref{tab:scenarios} lists the evaluations with their group and risk domain or fidelity level; Table~\ref{tab:provenance} gives the primary source of each one's facts, which is the reference cited for it: references 1 to \voiNScenarios{} of this report and of the paper are the evaluations, in id order. The selection was made before any result was seen (research/scenarios\_draft/SELECTION.md in the repository). LLM safety evaluations were scored by how many LLM system cards report them and how directly their hazard result can be checked (a programmatic or expert-baselined score over a judge validated against humans over a judge validated only by itself), with a penalty for capability benchmarks that do not measure a hazard; near-duplicates of one construct were passed over in favour of another risk domain, and every risk domain has at least three. Physical-AI safety evaluations were scored by how directly they measure a physically harmful event and by whether they publish validity evidence (judge agreement with humans, reproduction on hardware), which the elicitors need to estimate $s$ and $t$ from, and chosen to span the fidelity levels of the sim-to-real gap; the best evaluation at each level was kept. Every decision and instrument fact cites a source, and the source texts are cached in the repository.

\begin{table}[!htbp]
  \centering
  \caption{The evaluations. Domain for LLM evaluations; fidelity level (\voiLevelMin{} to \voiLevelMax) for physical-AI evaluations.}
  \label{tab:scenarios}
  \input{../generated/tab_scenarios.tex}
\end{table}

\begin{table}[!htbp]
  \centering
  \caption{Provenance: the primary source of each evaluation's facts, the reference the paper cites for it (the first source of its scenario file that is not a system card), as titled in the scenario file; ``and system cards'' marks evaluations whose facts also quote LLM system cards.}
  \label{tab:provenance}
  \resizebox{\linewidth}{!}{\input{../generated/tab_provenance.tex}}
\end{table}

%===============================================================================
\FloatBarrier
\section{The ensemble}
\label{sec:final}

The run uses six members from six developers through one routing service, one call per member and prompt, $B$ and $K$ valued for society, no numeric anchors and the curated facts. Three members come from Chinese developers and three from US developers. Members were chosen for developer diversity, reliable JSON output and price. Some evaluations come from member developers (Google DeepMind's ASIMOV and Veo evaluations), and two LLM evaluations from Anthropic, which has no member; we treat this as a limitation, not a correction. Each member is pinned to one endpoint with fallbacks disabled, so the weights and the price cannot change during the run; endpoints were chosen among those that keep no data. Reasoning effort is set to medium for every member. Requests ask for a JSON object and refuse endpoints that ignore the request options.

\begin{table}[!htbp]
  \centering\footnotesize
  \caption{The members (protocol \texttt{final}).}
  \label{tab:final}
  \begin{tabular}{@{}llll@{}}
    \toprule
    model & developer & endpoint & reasoning effort\\
    \midrule
    \texttt{deepseek/deepseek-v4.1-flash} & DeepSeek & DeepInfra (fp8) & medium\\
    \texttt{z-ai/glm-5.3} & Z.ai & Z.ai (fp8) & medium\\
    \texttt{xiaomi/mimo-v2.6-flash} & Xiaomi & DeepInfra (fp8) & medium\\
    \texttt{openai/gpt-6-luna} & OpenAI & Azure & medium\\
    \texttt{google/gemini-3.8-flash} & Google & Vertex & medium\\
    \texttt{x-ai/grok-4.7} & xAI & xAI & medium\\
    \bottomrule
  \end{tabular}
\end{table}

Why many members and no repeats. In the pilots the variance between members was comparable to the variance between repeats of one member, even with members from one developer, so a call spent on a new member buys at least as much as a call spent on a repeat. A diverse crowd of LLMs forecasts about as well as a human crowd~\citep{schoenegger2024wisdom}, and models of different developers make less correlated errors than models of one developer~\citep{kim2025correlated}. The counter-argument, that a few strong models beat a mix, assumes a verifier that tells which model is best; for elicited priors there is none.

%===============================================================================
\FloatBarrier
\section{The most damaging things a reviewer could say}
\label{sec:reviewer}

Each objection is stated as a reviewer would put it, followed by our answer.

\begin{enumerate}
\item \emph{``Language models made up the numbers.''} They did, and the paper says so. The claim is comparative: where one group sits relative to the other, under one set of beliefs applied to both. Every member answers every evaluation from the same template and the same kind of sourced facts, so a bias common to all evaluations moves both groups. A bias specific to physical-AI knowledge would not cancel, and the Monte Carlo interval does not capture it. Where a measured sensitivity, specificity or cost exists, it can replace the elicited one without changing the method.
\item \emph{``The members are one population. Their agreement proves nothing.''} Agreement between models trained on similar data cannot separate shared knowledge from shared bias~\citep{kim2025correlated}. The six members come from six developers, which is as far as diversity goes today. We report cross-member agreement only as evidence about the inputs (Fig.~\ref{fig:members}).
\item \emph{``You wrote the decisions, so you rank the decisions, not the evaluations.''} Partly true, and the paper's limitations say so. Most of the spread of $\etast$ comes from the stakes of the decision attached to each evaluation. Each decision is the release decision the evaluation's own sources describe, written to a fixed schema, and the decision prompt never sees the evaluation. The question asked is the one a developer faces: what is this evaluation worth for the decision it is used for. Pricing physical-AI and LLM evaluations against one shared decision, matched pairs, would separate the two.
\item \emph{``A binary pass or fail is a caricature of an evaluation that returns a score.''} The model is simple and aims at orders of magnitude and rankings. A score read against a threshold set for the decision can only widen the range of priors for which the result changes the decision. $\etast$, which does not depend on the threshold, places the groups the same way. % CHECK-DATA
\item \emph{``Most of your physical-AI evaluations test a language or vision-language model inside a robot, and one is a car.''} That is what published physical-AI safety evaluation looks like today, and it cuts both ways: the evaluations are few and imperfect, so the result says as much about the state of the field as about physical AI. The set spans fidelity levels from text-only planning to track testing, and the fidelity analysis (Fig.~\ref{fig:level}) uses that spread.
\item \emph{``The groups are small and chosen by different rules.''} The two lists are ranked by fixed rules written down before any result, and the test is exact. It separates the groups despite low power, and the separation has a named cause, the stakes per decision, that a reader can check against Figure~\ref{fig:params}. % CHECK-DATA
\item \emph{``Society's stakes are arbitrary.''} They are the stakes a published safety framework asks a developer to weigh. A developer's own exposure is smaller and would shrink both groups; the ratio between them is what the comparison rests on.
\item \emph{``Many evaluations have zero value, so the ranking is mostly ties.''} Under $\eta$, \voiZeroSharePhysical{} of the physical-AI and \voiZeroShareLLM{} of the LLM evaluations are worth close to zero under our model at the median parameters, and the paper reports that share as a finding: for those decisions the developer should mitigate first and evaluate the mitigated system. $\etast$ has no zeros and gives a full ranking. % CHECK-DATA
\item \emph{``Drawing parameters independently distorts the uncertainty.''} It does, in a direction we cannot sign: EVSI is not monotone in $p$, $B$ and $K$, and independence across evaluations narrows the interval of $\PS$. Drawing whole answers instead of parameters would measure the effect.
\item \emph{``Models refuse biosecurity and cyber prompts, so the sample is biased.''} Refusals are counted per member and class (Table~\ref{tab:health}); a slot with no valid answer after one retry is missing, not imputed.
\item \emph{``The context is cherry-picked.''} The facts are sentences an agent selected from the sources, so selection is possible. Every fact cites a source whose text is cached in the repository, the same facts go to every member, and the no-context ablation (Section~\ref{sec:noctx}) removes the selection altogether.
\end{enumerate}

%===============================================================================
\FloatBarrier
\section{What the pilots found}
\label{sec:pilots}

Two pilot studies preceded this one, with fewer evaluations and members from a single developer. The first priced ten LLM and five physical-AI evaluations. Its ranking by $\etast$ was stable across members, but mostly because the members agreed on the stakes of the decisions; whether a single result could change the decision varied between members. With five physical-AI evaluations it could not place the groups. The second priced evaluations of increasing fidelity for two fixed decisions. A single prompt let the instrument description move the elicited prior; separate decision and instrument prompts removed that effect. Its $\etast$ fell with fidelity level on both decisions. This study adopts both lessons: separate prompts, more evaluations per group, members from several developers, and fidelity as one discussion point. An iteration run with two Claude models (archived in the repository) preceded the final ensemble and drove the choice of evaluations and prompts; its numbers are not reported here because its templates still elicited a reuse count that the final model dropped.

%===============================================================================
\FloatBarrier
\section{Future work}
\label{sec:future}

Three extensions follow directly from the model. \emph{Growth over time}: as robot fleets grow, the stakes of a physical-AI release grow with them, and at a fixed threshold $\EVSI$ and $\EVSIst$ grow in proportion. \emph{Matched decisions}: pricing a physical-AI and an LLM evaluation against one shared decision separates the evaluation from the decision it informs. \emph{Other decision-makers}: a regulator permitting deployment, a funder building an evaluation for the field, and an insurer setting a premium. The project's future-work notes (docs/FUTURE\_WORK.md in the repository) list these in more detail.

\bibliography{../refs}

\end{document}
